\unit{Exposé 182 — Formal Geometry and Algebraic Geometry}
\sid{182}
\src{193}
\seminaren{(May 1959)}
\paperhead{FORMAL GEOMETRY AND ALGEBRAIC GEOMETRY}%
  {by Alexander GROTHENDIECK}

\fsection{1}{Schemes}

An affine algebraic space defined over a field $k$ is known to be essentially
determined by its affine algebra $A$ (the ring of regular functions defined over
$k$), with morphisms of algebraic spaces $X\longrightarrow Y$ corresponding
bijectively to homomorphisms of $k$-algebras
$A(Y)\longrightarrow A(X)$. The affine algebra corresponding to an algebraic
space is a finitely generated $k$-algebra and, from the “classical” point of
view, has no nilpotent elements; conversely, every algebra of this type is
obtained as the affine algebra of an algebraic space defined over $k$. There is
then a familiar dictionary for interpreting situations concerning affine
algebraic spaces in terms of commutative algebra. It has long been observed
that this yields more general statements, since it is usually no longer
necessary to suppose that the rings in question are of the type just
considered, the Noetherian hypothesis most often being sufficient. In
particular, whether or not a base field is specified, there is no reason to
exclude the case in which these rings contain nilpotent elements. Until now,
geometers have refused to take these indications into account and have
persisted in restricting themselves to affine algebras without nilpotent
elements, i.e. to algebraic spaces whose structure sheaves contain no
nilpotent elements (and, most often, even to “absolutely irreducible”
algebraic spaces). The speaker believes that this state of mind has been a
serious obstacle to the development of genuinely natural methods in algebraic
geometry.

Let $A$ be a commutative ring. It is well known that the set
$X=\operatorname{Spec}(A)$ of prime ideals of $A$ carries a natural topology,
the “Zariski topology” or spectral topology. Moreover, there is a sheaf of
commutative rings $\Osh_X$ on $X$ whose stalk at $\mathfrak p\in X$ is the
localized ring $A_{\mathfrak p}$ and whose ring of sections identifies with
$A$. Thus $X$ becomes a \emph{ringed space}, called the \emph{prime spectrum}
of $A$. A ring homomorphism $f:A\longrightarrow B$ defines a morphism of
ringed spaces
$f':\operatorname{Spec}(B)\longrightarrow\operatorname{Spec}(A)$ whose
underlying map of sets is simply
$\mathfrak p\longmapsto f^{-1}(\mathfrak p)$. The homomorphisms

\src{194}
of ringed spaces from $\operatorname{Spec}(B)$ to
$\operatorname{Spec}(A)$ obtained in this way are exactly those for which the
homomorphisms $\Osh_x\longrightarrow\Osh_y$ ($x=f'(y)$) are local (i.e. the
inverse image of the maximal ideal is the maximal ideal).

An \emph{affine scheme} is a ringed space isomorphic to some
$\operatorname{Spec}(A)$, and a \emph{prescheme} is a locally affine ringed
space, i.e. one in which every point has an open neighborhood that is an affine
scheme for the induced structure. Morphisms of preschemes are defined in the
evident way; locally they correspond to ring homomorphisms.

When a prescheme $S$ is fixed and morphisms of preschemes
$X\longrightarrow S$ are considered, $S$ plays the role of a base field or
base ring (or better, of a base space in a fibration). One then says that $X$
is an $S$-prescheme; if $S=\operatorname{Spec}(A)$, this also means that
$\Osh_X$ is a sheaf of $A$-algebras. Thus every prescheme may be regarded
uniquely as a $\mathbf Z$-prescheme. The $S$-preschemes naturally form a
category; moreover, one proves that the product of any two objects $X,Y$ in
this category always exists, and it is denoted $X\times_S Y$. This notion of
product makes it possible to define base change in an $S$-prescheme
corresponding to a morphism $S'\longrightarrow S$: indeed, $X\times_S S'$ may
be regarded as an $S'$-prescheme.

One says that $X$ is \emph{separated} over $S$ if the diagonal of
$X\times_S X$ is closed. A \emph{scheme} is a prescheme separated over
$\mathbf Z$; it is then separated over anything. For simplicity, from now on
we shall speak only of schemes, which we shall moreover assume
\emph{Noetherian}, i.e. finite unions of affine open subsets that are spectra
of Noetherian rings. One says that $X$ is \emph{of finite type over} $S$ if,
for every affine open subset $U$ of $S$, its inverse image in $X$ is a finite
union of affine open subsets whose rings are finite-type algebras over the ring
of $U$. It is such $S$-schemes that lend themselves to a properly geometric
study. In particular, for every $s\in S$, the fiber $f^{-1}(s)$ of $X$ over
$s$ is an algebraic scheme over the residue field $\kappa(s)$ of the local ring
$\Osh_s$ of $s$ in $S$. Thus $X$ may, to a certain extent, be regarded as a
family of “algebraic spaces” $f^{-1}(s)$, with the parameter $s$ ranging over
$S$ (i.e., from the local point of view, over the set of prime ideals of a
given ring). The $\kappa(s)$ may of course have different characteristics. If
$S=\operatorname{Spec}(k)$, where $k$ is a field, one essentially recovers the
usual notion of an “algebraic space,” with the sole difference that the
structure sheaf may now have nilpotent elements.

\src{195}
Drawing on familiar notions, one defines a \emph{projective morphism}, and more
generally a \emph{proper morphism}. Such a morphism is of finite type; moreover,
it maps closed subsets to closed subsets and retains this property under
arbitrary base change.

If $X$ is a scheme (Noetherian, as always), the sheaf $\Osh_X$ is a
\emph{coherent sheaf of rings} in the sense of [2]. Thus the coherent sheaves
of modules on $X$ are also precisely those sheaves locally isomorphic to the
cokernel of a homomorphism $\Osh_X^{m}\longrightarrow\Osh_X^{n}$.

\fsection{2}{Formal schemes}

Let $X$ be a scheme and $X'$ a closed subset of $X$. There is then a coherent
subsheaf $J$ of $\Osh_X$ such that
$X'=\operatorname{Supp}(\Osh_X/J)$ (and there is even a largest such
subsheaf). Equipped with $\Osh_X/J$, $X'$ becomes a scheme, denoted $X_0$; such
a scheme is called a “closed subscheme of $X$.” For every $n$, one may also
consider $X'$ equipped with $\Osh_X/J^{n+1}$, denoted $X_n$; this is a closed
subprescheme of $X$ whose underlying set is still $X'$ but which has a
different structure sheaf, namely $\Osh_{X_n}=\Osh_X/J^{n+1}$. The
$\Osh_{X_n}$ plainly form an inverse system of sheaves of rings on $X'$, whose
inverse limit $\overline{\Osh}_X$ is called the \emph{formal completion of
$\Osh_X$ along $X'$}. Equipped with this sheaf of rings, $X'$ is called the
\emph{formal completion of $X$ along $X'$}; it is therefore a ringed space but
not, in general, a scheme. For every coherent sheaf $F$ on $X$, one may
similarly consider the formal completion
$\overline F=\varprojlim_n F_n$ of $F$ along $X'$ (where
$F_n=F\otimes_{\Osh_X}\Osh_X/J^{n+1}$); it is a sheaf of modules on
$\overline X$. Its sections are called the “formal sections of $F$ along
$X'$” and identify with the elements of
$\varprojlim_n\Gamma(X',F_n)$. For $F=\Osh_X$, one obtains the “holomorphic
functions” of $X$ along $X'$ in Zariski's sense (whose terminology we shall
not follow because it conflicts with classical terminology).

A \emph{formal scheme} (understood to mean Noetherian) is a topological space
$\mathfrak X$ equipped with a sheaf of topological rings
$\Osh_{\mathfrak X}$ satisfying the following condition: there is an
isomorphism of sheaves of topological rings
$\Osh_{\mathfrak X}=\varprojlim_n\Osh_n$, where the $\Osh_n$ form an inverse
system of sheaves of rings on $\mathfrak X$, each making $\mathfrak X$ into a
scheme $\mathfrak X_n$, and such that for $m\geq n$ the homomorphism
$\Osh_m\longrightarrow\Osh_n$ is surjective and has kernel $J_m^{n+1}$, where
$J_m$ is

\src{196}
the kernel of $\Osh_m\longrightarrow\Osh_0$. One then proves that
$\Osh_{\mathfrak X}$ is a \emph{coherent sheaf of Noetherian local rings}.

By the definitions, a formal completion $\overline X$ as above is a formal
scheme, and conversely every formal scheme is \emph{locally} of this type. In
fact, specifying an \emph{affine} formal scheme (i.e. one for which
$\mathfrak X_0$ is affine, which implies that all the $\mathfrak X_n$ are
affine) is equivalent to specifying a Noetherian topological ring that is
separated and complete for the $J$-adic topology.

The usual definitions for ordinary schemes---morphism, morphism of finite type,
proper morphism, and so forth---extend without difficulty to formal schemes.

\fsection{3}{The three fundamental theorems}

Let $f:X\longrightarrow Y$ be a proper morphism of schemes (Noetherian, as
always), let $Y'$ be a closed subset of $Y$, and let $X'$ be its inverse image
in $X$; consider the corresponding formal completions $\overline Y$ and
$\overline X$. Then $f$ induces a morphism
$\overline f:\overline X\longrightarrow\overline Y$ of formal schemes, which
is itself proper. If $F$ is a coherent sheaf on $X$, then $\overline F$ is a
coherent sheaf on $\overline X$. In Theorem 1, the original $X,Y,F$ are
forgotten and only the proper morphism $\overline f$ of formal schemes and the
coherent sheaf $\overline F$ on $\overline X$ are retained. (The speaker has,
however, written out a complete proof only in the case that begins with
$X,Y,f,F$.)

\result{THEOREM}{1 (finiteness theorem)}
\begin{enumerate}
  \item[\textit{i.}] The $R^{q}\overline f_{*}(\overline F)$ are coherent
    sheaves on $\overline Y$.
  \item[\textit{ii.}] The natural homomorphisms
    \[
      R^{q}\overline f_{*}(\overline F)
      \longrightarrow \varprojlim_n R^{q}f_{n*}(F_n)
    \]
    are isomorphisms.
\end{enumerate}

In this statement, a coherent subsheaf $J$ of $\Osh_Y$ defining $Y'$ is
supposed chosen; its inverse image gives a coherent subsheaf of $\Osh_X$
defining $X'$, and hence the definition of $F_n,X_n,Y_n$ and
$f_n:X_n\longrightarrow Y_n$ as in Section 2. The minor changes needed to make
these notations explicit are evident if one starts with an arbitrary proper
morphism of two formal schemes.

Theorem 1 concerned only “formal cohomology.” The following theorem

\src{197}
relates it to “algebraic cohomology” and is akin to a well-known theorem of
Serre [4] comparing algebraic and analytic cohomology.

\result{THEOREM}{2 (First comparison theorem)} The $R^{q}f_{*}(F)$ are
coherent sheaves on $Y$ (which is a special case of Theorem 1), and the natural
homomorphisms
\[
  \overline{R^{q}f_{*}(F)}
  \longrightarrow \varprojlim_n R^{q}f_{n*}(F_n)
\]
are isomorphisms.

\result{COROLLARY}{1} There are canonical isomorphisms
\[
  \overline{R^{q}f_{*}(F)}=R^{q}\overline f_{*}(\overline F).
\]

For $q=0$, this corollary is a generalization of Zariski's “fundamental theorem
of holomorphic functions,” from which we shall deduce a generalization of
Zariski's “connectedness theorem.” Note, moreover, that although Theorem 1
(\textit{ii}) is trivial for $q=0$, the same is by no means true of Theorem 2
or of its equivalent formulation (Corollary 1). In fact, the proof proceeds by
descending induction on $q$ (being trivial for large $q$, since both sides then
vanish), and the case $q=0$ thus appears as the last step of the induction and
therefore, so to speak, as the “most difficult” case.

\result{COROLLARY}{2} Suppose $Y=\operatorname{Spec}(A)$, with $Y'$ defined by
the ideal $J$ of $A$. Then, for every coherent sheaf $F$ on $X$, the
$H^{q}(X,F)$ are finite-type $A$-modules whose $J$-adic completions are the
$H^{q}(\overline X,\overline F)$.

Finally, applying this corollary to
$H=\Homsh_{\Osh_X}(F,G)$ gives:

\result{COROLLARY}{3} Suppose $Y=\operatorname{Spec}(A)$, with $Y'$ defined by
the ideal $J$ of $A$. If $F,G$ are two coherent sheaves on $X$, then
$\operatorname{Hom}(F,G)$ is a finite-type $A$-module whose $J$-adic
completion identifies with
$\operatorname{Hom}(\overline F,\overline G)$.

The natural map
$\operatorname{Hom}(F,G)\longrightarrow
\operatorname{Hom}(\overline F,\overline G)$ is of course the one that assigns
to a homomorphism $u:F\longrightarrow G$ its extension “by continuity”
$\overline u:\overline F\longrightarrow\overline G$ (with which convention
$\overline F$ becomes a functor of $F$).

Now suppose that $A$ is separated and complete for its $J$-adic topology. The
preceding Corollaries 2 and 3 then give
\[
  H^{q}(X,F)=H^{q}(\overline X,\overline F),
  \qquad
  \operatorname{Hom}(F,G)
  =\operatorname{Hom}(\overline F,\overline G).
\]

\src{198}
This last identity shows that the category of coherent sheaves on $X$
identifies with a \emph{subcategory} (whose morphisms are the induced
morphisms) of the category of coherent sheaves on $\overline X$. In fact, one
has the stronger statement:

\result{THEOREM}{3} A sheaf of modules on $\overline X$ is coherent if and only
if it is isomorphic to a sheaf of the form $\overline F$, where $F$ is a
coherent sheaf on $X$ (determined up to canonical isomorphism by Theorem 2,
Corollary 3). [Recall that now $Y=\operatorname{Spec}(A)$, where $A$ is a
Noetherian topological ring that is separated and complete for the $J$-adic
topology.]

\result{COROLLARY}{1} The closed subschemes of $X$ correspond bijectively to
the closed formal subschemes of $\overline X$.

Indeed, they correspond to the coherent subsheaves of $\Osh_X$, respectively
of $\overline{\Osh}_X$. Viewing the graphs of morphisms as closed subschemes,
one deduces from Corollary 1:

\result{COROLLARY}{2} Let $X,Z$ be two schemes proper over $A$ (a Noetherian
$J$-adic ring that is separated and complete). Then the map
$g\longmapsto\overline g$ defines a bijection between the $Y$-morphisms from
$X$ to $Z$ and the $\overline Y$-morphisms from $\overline X$ to
$\overline Z$.

In other words, algebraic schemes proper over $A$ appear as a
\emph{subcategory} (whose morphisms are the induced morphisms) of the category
of formal schemes proper over $\overline Y$. It should be noted, however, that
there are formal schemes proper over $\overline Y$ that are not
“algebraizable,” i.e. are not isomorphic to any $\overline X$ with $X$ proper
over $A$ (just as there are compact complex analytic varieties that do not
come from algebraic varieties defined over the complex numbers). Such formal
schemes arise naturally in the “theory of moduli.” Let us nevertheless note
an interesting special case in which a formal scheme is algebraizable:

\result{THEOREM}{4} Let $A$ be a complete Noetherian local ring with residue
field $k$, and let $\mathfrak X$ be a formal scheme proper over $A$ (equipped
with its $\mathfrak m(A)$-adic topology). Suppose:
\begin{enumerate}
  \item[\textit{i.}] that the local rings of $\Osh_{\mathfrak X}$ are flat
    $A$-modules, or equivalently, if $\Osh_{\mathfrak X}$ and $A$ are equipped
    with the filtrations defined by the powers of the maximal ideal of $A$,
    that the associated graded objects satisfy
    \[
      \operatorname{gr}(\Osh_{\mathfrak X})
      \simeq
      \operatorname{gr}^{0}(\Osh_{\mathfrak X})
      \otimes_k\operatorname{gr}(A).
    \]
  \item[\textit{ii.}] regarding
    $\mathfrak X_0=\mathfrak X\otimes_A k$ as an algebraic scheme over $k$, one
    has
    \[
      H^2(\mathfrak X_0,\Osh_{\mathfrak X_0})=0.
    \]
  \item[\textit{iii.}] $\mathfrak X_0$ is projective.
\end{enumerate}
Under these conditions, $\mathfrak X$ is algebraizable; more precisely, it is
isomorphic to $\overline X$, where $X$ is a projective $A$-scheme.

Conditions \textit{(ii)} and \textit{(iii)} hold in particular when
$\mathfrak X_0$ is a nonsingular curve over $k$, and Theorem 4 can therefore
be applied, in particular, in the “theory of moduli” for curves of given
genus\ldots{} Let us indicate how Theorem 4 is proved. One shows (cf.
Proposition 3 below) that \textit{(i)} and \textit{(ii)} imply that every
coherent sheaf on $\mathfrak X_0$ which is locally isomorphic to the structure
sheaf is obtained by reduction from a sheaf of the same kind on
$\mathfrak X$. Starting with an “ample” sheaf on $\mathfrak X_0$ (one exists
by \textit{(iii)}), one then lifts it to an invertible sheaf on $\mathfrak X$;
using Theorem 1, one proves that a multiple of this sheaf defines an immersion
of $\mathfrak X$ into the formal completion of a scheme $\mathbf P_A^r$ (of
“projective type” and dimension $r$ over $A$).

For the proofs of Theorems 1 through 3, see [1].

\src{199}
\fsection{4}{Application to the connectedness theorem and Zariski's Main
Theorem}

Let $f:X\longrightarrow Y$ be a proper morphism of schemes. By the finiteness
theorem, $f_*(\Osh_X)=\mathcal A$ is a coherent sheaf on $Y$; it is moreover
a sheaf of commutative algebras, and therefore corresponds to a $Y$-scheme
$g:Y'\longrightarrow Y$ finite over $Y$ (defined by the conditions that it be
affine over $Y$, i.e. that the inverse image of an affine open subset be
affine, and that $g_*(\Osh_{Y'})=\mathcal A$). It is immediate that $f$ then
factors canonically as $f=gf'$, where $f':X\longrightarrow Y'$ is a morphism
from $X$ to $Y'$ satisfying $f'_*(\Osh_X)=\Osh_{Y'}$. This factorization of
$f$ is called the \emph{Stein factorization} of $f$. Applying the first
comparison theorem, Corollary 1, to $f'$ and to the part of $Y'$ reduced to a
point $y'$, one finds that $f'^{-1}(y')=X'$ is connected (otherwise the formal
sections of $X$ along $X'$ would not form a local ring, whereas the completion
of $f'_*(\Osh_X)_{y'}=\Osh_{Y',y'}$ is local!). We have proved:

\result{THEOREM}{5 (Zariski's “connectedness theorem”)} Let
$f:X\longrightarrow Y$ be a proper morphism. Then $f$ factors uniquely (up to
isomorphism) as $f=gf'$, where $g:Y'\longrightarrow Y$ is finite and
$f':X\longrightarrow Y'$ satisfies $f'_*(\Osh_X)=\Osh_{Y'}$ (whence
$g_*(\Osh_{Y'})=f_*(\Osh_X)$). The fibers of $f'$ are connected; that is, the
set of connected components of a fiber $f^{-1}(y)$ of $f$ is in bijection
with the set of points of $Y'$ above $y$, i.e. with the set of maximal ideals
of $f_*(\Osh_X)_y$.

The usual variants of the connectedness theorem follow immediately. We note
here only the following.

\src{200}
\result{COROLLARY}{1} A point $x$ of $X$ is isolated in its fiber
$f^{-1}(y)$ if and only if the fiber $f'^{-1}(y')$ (where $y'=f'(x)$) is
reduced to $x$, or equivalently if $f'$ induces an isomorphism from a
neighborhood of $x$ onto a neighborhood of $y'$. The set of these points is
an open subset $U$, and $f'$ induces an isomorphism from $U$ onto an open
subset of $Y'$.

To show that $f'$ is a local isomorphism at $x$, observe that $f'$ induces an
isomorphism $\Osh_{y'}\longrightarrow\Osh_x$, as follows from
$f'_*(\Osh_X)=\Osh_{Y'}$ and from the fact that (since $f'$ is a closed map
whose fiber over $y'$ is reduced to $x$) the sets $f'^{-1}(V)$ run through a
fundamental system of neighborhoods of $x$ as $V$ runs through a fundamental
system of neighborhoods of $y'$. The following result, due to Chevalley in the
“geometric” case, follows at once.

\result{COROLLARY}{2} A morphism $f$ is finite if and only if it is proper and
its fibers are finite.

Indeed, in that case $f'$ is an isomorphism by the preceding result.

Let $f:X\longrightarrow Y$ be a morphism not necessarily proper, but suppose
that $X$ is contained as an open subset in a proper $Y$-scheme
$\widetilde f:\widetilde X\longrightarrow Y$ (as will be the case, in
particular, if $f$ is quasi-projective). Applying Corollary 1, one sees that
$\widetilde f'$ induces an isomorphism from the set $U$ of points of $X$
isolated in their fiber onto an open subset of $Y'$ (and that $U$ is indeed
open). This gives the following global form of Zariski's Main Theorem.

\result{THEOREM}{6 (Zariski's Main Theorem)} Let
$f:X\longrightarrow Y$ be a morphism of finite type. Then the set $U$ of
points of $X$ that are isolated in their fiber is open; and if $f$ is
separated,\footnote{The 1962 erratum replaces the printed hypothesis “$f$ is
quasi-projective” by the weaker hypothesis “$f$ is separated,” by SGA VIII,
6.2: every separated quasi-finite morphism is quasi-projective.} $U$ is
$Y$-isomorphic to an open part of a scheme $Y'$ finite over $Y$.

Since a morphism of finite type is locally affine and hence,
\textit{a fortiori}, locally quasi-projective, the usual local forms of the
Main Theorem follow at once from Theorem 6.

\fsection{5}{Application to the cohomological study of proper flat morphisms}

Let $f:X\longrightarrow Y$ be a proper morphism, and let $F$ be a coherent
sheaf on $X$, assumed flat over $Y$, i.e. the $F_x$ are flat modules over the
rings $\Osh_y$ ($y=f(x)$). Equivalently, for every $y\in Y$, if $F$ is
filtered along the fiber $f^{-1}(y)$ by the $\mathfrak m_y^nF$ (where
$\mathfrak m_y$ is the maximal ideal of $\Osh_y$), the associated graded
object is isomorphic to
$(F/\mathfrak m_yF)\otimes_{\kappa(y)}\operatorname{gr}(\Osh_y)$; in other
words,
\[
  \mathfrak m_y^nF/\mathfrak m_y^{n+1}F
  =F_y\otimes_{\kappa(y)}
  (\mathfrak m_y^n/\mathfrak m_y^{n+1})
\]
for every integer $n$. Here $X_y$ denotes the fiber $f^{-1}(y)$ regarded as a
proper scheme over the residue field $\kappa(y)$ of $y$, and $F_y$ is the
sheaf $F/\mathfrak m_yF$ induced by $F$ on $X_y$. Taking this isomorphism and
Theorem 2 into account, one obtains bounds, and sometimes computations, for
the $R^qf_*(F)$ near $y$ from the cohomology of $X_y$ with coefficients in
$F_y$. Here Theorem 2 takes the form
\[
  \widehat{R^qf_*(F)_y}
  =\varprojlim_n H^q(X_y,F/\mathfrak m_y^nF).
\]
We shall mention only the following consequence.

\src{201}
\result{PROPOSITION}{1} Let $f:X\longrightarrow Y$ be a proper morphism and
let $F$ be a coherent sheaf on $X$, flat over $Y$. Let $y\in Y$ and let $q$
be an integer, and suppose that $H^q(X_y,F_y)=0$. Then $R^qf_*(F)$ vanishes in
a neighborhood of $y$, and for every $n$ the natural homomorphism
\[
  R^{q-1}f_*(F)_y\longrightarrow
  H^{q-1}(X_y,F/\mathfrak m_y^nF)
\]
is surjective.

In particular, if $f$ is a flat morphism (i.e. $\Osh_X$ is flat over $Y$),
then every locally free coherent sheaf $F$ on $X$ is flat over $Y$. Let $F,G$
be two such sheaves. Applying Proposition 1 to
$\Homsh_{\Osh_X}(F,G)$ with $q=1$, one obtains:

\result{THEOREM}{7} Let $f$ be a proper flat morphism, let $F,G$ be two
locally free coherent sheaves on $X$, let $y\in Y$, and suppose that
\[
  H^1\bigl(X_y,\Homsh_{\Osh_{X_y}}(F_y,G_y)\bigr)=0.
\]
Then every homomorphism $u_0:F_y\longrightarrow G_y$ is induced by a
homomorphism
\[
  u:F|V\longrightarrow G|V,
  \qquad V=f^{-1}(U),
\]
where $U$ is a neighborhood of $y$.

\result{COROLLARY}{} If $u_0$ is an isomorphism (respectively a monomorphism,
an epimorphism), then so is $u$ for sufficiently small $U$.

In particular:

\src{202}
\result{COROLLARY}{2} Let $E_0$ be a locally free coherent sheaf on $X_y$ such
that
\[
  H^1\bigl(X_y;\Homsh_{\Osh_{X_y}}(E_0,E_0)\bigr)=0.
\]
Then two locally free sheaves whose restrictions to $X_y$ are isomorphic to
$E_0$ are isomorphic in a neighborhood of $X_y$.

Thus:

\result{COROLLARY}{3} Suppose that $H^1(X_y,\Osh_{X_y})=0$. Then two
invertible sheaves on $X$ (i.e. sheaves locally isomorphic to $\Osh_X$) whose
restrictions to $X_y$ are isomorphic are themselves isomorphic in a neighborhood of $X_y$.\footnote{The printed text does not specify the neighborhood. The local scope of Corollary 2 is restored here; it is also used in the proof of Proposition 2.}

It follows that:

\result{PROPOSITION}{2} Let $Y$ be a connected scheme and let $E$ be a locally free coherent sheaf of rank at least $2$ on $Y$. Consider the associated projective-space bundle
$X=\mathbf P(E)$, equipped with its familiar invertible sheaf $\Osh_X(1)$.
Then every invertible sheaf $L$ on $X$ is isomorphic to a sheaf of the form
\[
  f^*(L')\otimes\Osh_X(n),
\]
where $L'$ is an invertible sheaf on $Y$ and $n$ is an integer. The integer
$n$ is uniquely determined, and $L'$ is determined up to isomorphism.

The preceding Corollary 3 proves that $L$ is isomorphic to some $\Osh_X(n)$ in
a neighborhood of every fiber. The rest is more or less formal.

Proposition 2 makes it possible to determine the $Y$-morphisms from
$X=\mathbf P(E)$ to another projective bundle. In particular, one finds:

\result{COROLLARY}{} Let $u$ be an automorphism of $X=\mathbf P(E)$. Then
there exist an invertible sheaf $L'$ on $Y$ and an isomorphism
$v:E\longrightarrow E\otimes L'$ such that $u$ is the corresponding
isomorphism
\[
  \mathbf P(E)\xrightarrow{\sim}\mathbf P(E\otimes L')=\mathbf P(E);
\]
the pair $(v,L')$ is determined up to isomorphism.

Let $\Gamma$ be the set of classes of invertible bundles $L'$ on $Y$ such that
$E\otimes L'$ is isomorphic to $E$. Its elements are torsion: indeed, if $n$
is the rank of $E$, taking $n$th exterior powers shows that one must have
$L'^{\otimes n}\simeq\Osh_Y$. The corollary may therefore be expressed by
saying that there is an exact sequence of groups
\[
  e\longrightarrow
  \operatorname{Aut}(E)/\Gamma(Y,\Osh_Y^*)
  \longrightarrow \operatorname{Aut}_Y(X)
  \longrightarrow \Gamma\longrightarrow e
\]
(which may also be deduced from the cohomology exact sequence arising from the
following exact sequence of sheaves of groups
\[
  e\longrightarrow \Osh_Y^*
  \longrightarrow \underline{\operatorname{Aut}}(E)
  \longrightarrow \underline{\operatorname{Aut}}_Y(X)
  \longrightarrow e,
\]
where $\Osh_Y^*$ is the sheaf of “units” of $\Osh_Y$, identified with the
center of $\underline{\operatorname{Aut}}(E)$).

\src{203}
\fsection{6}{Application to existence and uniqueness theorems for sheaves and
schemes over a complete $J$-adic ring}

Theorem 7 gave a uniqueness result for locally free coherent sheaves, using
Theorems 1 and 2. Using Theorem 3, we shall now deduce existence theorems for
sheaves, morphisms of schemes, or schemes. In what follows, $A$ denotes a
Noetherian local ring that is separated and complete. The general method
always consists in carrying out formal constructions, which amounts
essentially to doing algebraic geometry over an Artinian ring, and then using
the three fundamental theorems to draw conclusions of an “algebraic” nature.

\result{PROPOSITION}{3} Let $\mathfrak X$ be a formal scheme proper and flat
over $A$, and let $F_0$ be a locally free sheaf on $X_0$ such that
\[
  H^2\bigl(X_0,\Homsh_{\Osh_{X_0}}(F_0,F_0)\bigr)=0.
\]
Then there exists a locally free sheaf $F$ on $\mathfrak X$ whose restriction
to $X_0$ is isomorphic to $F_0$. (Moreover, this $F$ is unique up to
isomorphism if
$H^1(X_0,\Homsh_{\Osh_{X_0}}(F_0,F_0))=0$.)

One constructs, successively, locally free sheaves $F_n$ on the $X_n$, each
inducing the preceding one. The construction of $F_n$ encounters an
obstruction in
\[
  H^2\bigl(X_0,\Homsh_{\Osh_{X_0}}(F_0,F_0)\bigr)
  \otimes_{A/J}(J^n/J^{n+1}),
\]
which therefore vanishes by hypothesis. Applying Theorem 3 now gives:

\result{COROLLARY}{1} Let $X$ be a scheme proper and flat over $A$, and let
$F_0$ be as above. Then there exists a locally free sheaf $F$ on $X$ whose
restriction to $X_0$ is isomorphic to $F_0$. Moreover, $F$ is unique up to
isomorphism if $H^1(X_0,\Homsh_{\Osh_{X_0}}(F_0,F_0))=0$.

Let $X_0$ be a scheme of finite type over the field $k$. Suppose that $X_0$ is
nonsingular (by which we mean absolutely nonsingular) over $k$, but not
necessarily proper over $k$. Let $A$ be a local Artinian ring with residue
field $k$. We are interested in finding schemes $X$ flat over $A$ such that
$X\otimes_A k=X_0$ (this is the starting point of the “theory of moduli,” or
of “variations of structure” of $X_0$). Giving such an $X$ is equivalent to
giving, on the topological space $X_0$, a sheaf $\Osh_X$ endowed with the
following structures:
\begin{enumerate}
  \item[\textit{i.}] It is a sheaf of $A$-algebras.
  \item[\textit{ii.}] It is endowed with an augmentation homomorphism
    $\Osh_X\longrightarrow\Osh_{X_0}$ (compatible with the $A$-algebra
    structures), subject to the following conditions: the augmentation
    induces an isomorphism
    $\Osh_X\otimes_A k\xrightarrow{\sim}\Osh_{X_0}$; and $\Osh_X$ is flat
    over $A$, i.e. the graded object associated to $\Osh_X$, filtered by the
    powers of the maximal ideal $\mathfrak m$ of $A$, is isomorphic to
    $\operatorname{gr}^0(\Osh_X)\otimes_k\operatorname{gr}(A)$; equivalently,
    there are isomorphisms
    \[
      \mathfrak m^n\Osh_X/\mathfrak m^{n+1}\Osh_X
      =\Osh_{X_0}\otimes_k(\mathfrak m^n/\mathfrak m^{n+1}).
    \]
\end{enumerate}
The fundamental fact is the following.

\src{204}
\result{THEOREM}{8} Let $X_0$ be a scheme of finite type and nonsingular over
the field $k$, and suppose that $X_0$ is affine. Let $A$ be a local Artinian
ring with residue field $k$. Then there exists an $A$-scheme $X$, flat over
$A$, such that $X\otimes_A k=X_0$, and any two such schemes are necessarily
isomorphic.

The isomorphism in question is not canonical, since in general $X$ will have
nontrivial $A$-automorphisms inducing the identity on $X_0$. Moreover, in
general there is no “canonical” choice of an $X$ satisfying the stated
conditions, except when $A$ is a $k$-algebra (the equal-characteristic case),
where one may take $X=X_0\otimes_k A$, i.e.
$\Osh_X=\Osh_{X_0}\otimes_k A$ (whether or not $X_0$ is affine). In unequal
characteristic, when $X_0$ is not affine, I do not know in general whether
$X_0$ can be “lifted” to an $X$ defined over $A$. Nevertheless, let $n>0$ be
an integer, put $A_{n-1}=A/\mathfrak m^n$, and suppose that $X_0$ has been
lifted to a flat $A_{n-1}$-scheme $X_{n-1}$; we seek to lift $X_{n-1}$ to a
flat $A_n$-scheme $X_n$. Theorem 8 already shows that this is possible
locally. Moreover, one checks easily that if $U_n$ lifts an open subset
$U_{n-1}$ of $X_{n-1}$, then the sheaf of groups of automorphisms of $U_n$
(inducing the identity on $U_{n-1}$) is canonically isomorphic to
\[
  \Theta_{X_0/k}\otimes_k(\mathfrak m^n/\mathfrak m^{n+1})
\]
restricted to $U_{n-1}$, and in particular is commutative
($\Theta_{X_0/k}$ denotes the sheaf of germs of $k$-derivations on $X_0$).
It follows easily that there is an obstruction to constructing $X_n$ lifting
$X_{n-1}$, lying in
\[
  H^2(X_0,\Theta_{X_0/k})\otimes_k
  \mathfrak m^n/\mathfrak m^{n+1}.
\]
Consequently:

\result{COROLLARY}{1} Let $X_0$ be a scheme of finite type and nonsingular
over $k$, and suppose that $H^2(X_0,\Theta_{X_0/k})=0$. Then, for every local
Artinian ring $A$ with residue field $k$, there exists an $A$-scheme $X$, flat
over $A$, such that $X\otimes_A k=X_0$.

Moreover, if one has found an $X$ flat over $A$ lifting $X_0$, then by Theorem
8 the set of classes (up to isomorphism) of flat $A$-schemes lifting $X_0$
identifies with
\[
  H^1\bigl(X_0,\underline{\operatorname{Aut}}(X)\bigr),
\]
where $\underline{\operatorname{Aut}}(X)$ denotes, of course, the sheaf of
germs of automorphisms of the sheaf of $A$-algebras $\Osh_X$ that are
compatible with the augmentation. The filtration of $\Osh_X$ defines a
filtration of $\underline{\operatorname{Aut}}(X)$, the quotient of this sheaf
by the $n$th subgroup of the filtration being
$\underline{\operatorname{Aut}}(X_n)$. The graded object associated with this
filtration is commutative and identifies with
$\Theta_{X_0/k}\otimes_k\operatorname{gr}(A)$. In particular, if
$\mathfrak m^{n+1}$ is the first power of $\mathfrak m$ that vanishes, then
$F^n(\underline{\operatorname{Aut}}(X))$ (the last step of the filtration) is
contained in the center of $\underline{\operatorname{Aut}}(X)$ and is
isomorphic to $\Theta_{X_0/k}\otimes_k\mathfrak m^n$; it is also the sheaf of
germs of automorphisms of $X$ that induce the identity on
$X_{n-1}=X\otimes_A A/\mathfrak m^n$. These results immediately give the
following statements.

\src{205}
\result{COROLLARY}{2} Let $X_0$ be a scheme of finite type and nonsingular
over $k$. Let $A$ be a local Artinian ring with residue field $k$ and maximal
ideal $\mathfrak m$, suppose that $\mathfrak m^{n+1}=0$, put
$A_{n-1}=A/\mathfrak m^n$, and let $X_{n-1}$ be a flat $A_{n-1}$-scheme such
that $X_{n-1}\otimes_A k=X_0$. Then the set of classes (up to isomorphism
inducing the identity on $X_{n-1}$) of flat $A$-schemes $X_n$ such that
$X_n\otimes_A A_{n-1}=X_{n-1}$ is either empty or a principal homogeneous
space under
\[
  H^1(X_0,\Theta_{X_0/k})\otimes_k\mathfrak m^n.
\]

(Note that in general there is no privileged origin in this latter space,
since there is no privileged way to lift $X_{n-1}$ to $X_n$.)

\result{COROLLARY}{3} Let $X_0$ be a scheme of finite type and nonsingular
over $k$, and suppose that $H^1(X_0,\Theta_{X_0/k})=0$. Then, for every local
Artinian ring $A$ with residue field $k$, there exists at most one (up to
isomorphism) flat $A$-scheme $X$ such that $X\otimes_A k=X_0$.

Corollaries 1 and 3 immediately imply the apparently more general statements
obtained by assuming only that $A$ is a complete Noetherian local ring with
residue field $k$, provided that $X$ is introduced as a formal scheme over
$A$:

\result{THEOREM}{9} Let $k$ be a field and let $X_0$ be a scheme of finite
type and nonsingular over $k$. For every complete Noetherian local ring $A$
with residue field $k$, let $F(A)$ be the set of classes (up to isomorphism
inducing the identity on $X_0$) of formal schemes $\mathfrak X$ over $A$, of
finite type and flat over $A$, such that $\mathfrak X\otimes_A k=X_0$. With
this notation:
\begin{enumerate}
  \item[\textit{i.}] If $H^1(X_0,\Theta_{X_0/k})=0$, then for every $A$,
    $F(A)$ has at most one element.
  \item[\textit{ii.}] If $H^2(X_0,\Theta_{X_0/k})=0$, then for every $A$,
    $F(A)$ has at least one element.
\end{enumerate}

\src{206}
\result{COROLLARY}{1} Suppose that $X_0$ is proper over $k$. Under condition
\textit{(i)}, for every $A$ there exists at most one (up to isomorphism
inducing the identity on $X_0$) scheme $X$, proper and flat over $A$, such that
$X\otimes_A k=X_0$.

Here one uses Theorem 3, Corollary 2. For example:

\result{COROLLARY}{2} If $X$ is a scheme proper and flat over $A$ such that
$X\otimes_A k$ is isomorphic to the standard projective scheme
$\mathbf P_k^r$ of dimension $r$ over $k$, then $X$ is isomorphic to
$\mathbf P_A^r$.

(This result may also be deduced from Proposition 3, Corollary 1.)

\result{COROLLARY}{3} Let $X_0$ be a nonsingular projective scheme over $k$,
and suppose that
\[
  H^2(X_0,\Osh_{X_0})=H^2(X_0,\Theta_{X_0/k})=0.
\]
Then for every $A$ there exists a projective scheme $X$, flat over $A$, such
that $X\otimes_A k=X_0$.

Combine Theorem 9\textit{(ii)} with Theorem 4. In particular:

\result{COROLLARY}{4} Let $X_0$ be the scheme of a complete nonsingular
algebraic curve over $k$. Then, for every complete Noetherian local ring $A$
with residue field $k$, there exists a “smooth-curve scheme” $X$ over $A$ such
that $X\otimes_A k=X_0$.

\result{REMARKS}{}

1) Corollaries 3 and 4 are chiefly interesting when $k$ has characteristic
$p\ne0$, by taking for $A$ a discrete valuation ring of characteristic $0$
with residue field $k$; for example, the “smallest possible $A$,” namely one
whose maximal ideal is generated by $p$. (In fact, because of Cohen's
theorems, it is enough to know Corollaries 3 and 4 for one such ring $A$.) It
may be noted in this connection that, according to the specialists, it is not
known whether there exist schemes over a field $k$ that are not the reduction
modulo $p$ of a flat scheme defined over such a ring $A$. At least the results
of this section provide a method for systematically investigating this
question. One should begin
by examining whether the first obstruction in $H^2(X_0,\Theta_{X_0/k})$ is
necessarily zero.

\emph{1962 supplement.} --- J.-P. Serre constructed a nonsingular projective
threefold over an algebraically closed field $k$ of characteristic $p>0$ that
is not obtained by reduction from a proper scheme over an integral local ring
with residue field $k$ and fraction field of characteristic zero. Mumford is
said to have found an analogous example with a nonsingular projective surface.

2) It should be noted that Theorem 3, and the corresponding technique, applies
only when the base ring (local, to fix ideas) is complete.

To pass from results known for the completion of a local ring to corresponding
results for the local ring itself, one would need a fourth “fundamental
theorem,” whose definitive statement remains to be found.

\src{207}
3) The results of this section (especially the preceding Corollaries 1 and 2)
and of the next should be compared with the results of Kodaira--Spencer on the
variation of complex structures. Using the conjectural theorem just mentioned,
one should be able to conclude, under the hypotheses of Corollary 1 but without
assuming $A$ complete, that there exists a ring $A'$ containing $A$, finite and
unramified over $A$, such that $X\otimes_A A'$ and $X'\otimes_A A'$ are
$A'$-isomorphic (where $X,X'$ are two given proper flat $A$-schemes satisfying
$X\otimes_A k=X'\otimes_A k=X_0$). This can at least be proved when
$X_0=\mathbf P_k^r$, using the corollary to Proposition 2. In any case, when
$H^1(X_0,\Theta_{X_0/k})=0$, one can prove that the fibers of $X,X'$ over every
point $y$ of $Y=\operatorname{Spec}(A)$ are isomorphic, at least after passing
to the algebraic closure of the residue field $\kappa(y)$. (There is a local
result, apparently stronger, in which $A$ need no longer be assumed local.)
Concerning “variations of structure” for projective space, let us also mention
the following question, suggested by a corresponding problem of
Kodaira--Spencer. Let $X$ be a scheme proper and flat over an integral local
ring $A$ with fraction field $K$ and residue field $k$. Suppose that
$X\otimes_A K$ is isomorphic to $\mathbf P_K^r$. Is it true that
$X\otimes_A k=X_0$ is isomorphic to $\mathbf P_k^r$ (at least over the
algebraic closure of $k$)? In this problem one may assume that $A$ is a
complete discrete valuation ring. There is an analogous problem when $X_0$ is
an abelian variety.

\emph{1962 supplement.} --- The speaker is by then less optimistic about the
results conjectured above. The problem concerning variations of structure of
projective space was nevertheless solved affirmatively by Hironaka, and the
analogous problem for abelian varieties by Koizumi.

\fsection{7}{Application to the theory of moduli}

Since the speaker has only just begun to study this theory himself, we must
limit ourselves to a few indications. For simplicity we work over a field
$k$, i.e. in equal characteristic, although Theorem 8 also makes it possible
to treat the general case, apparently without essential change. At present we
have not gone beyond the “formal” stage; nevertheless, the speaker hopes to be
able to construct genuine moduli schemes in certain cases and, in particular,
for every integer $g$, to construct a scheme over the integers playing the
role of a universal moduli scheme for nonsingular curves of genus $g$.

\emph{1962 supplement.} --- Mumford constructed the moduli schemes for curves
of genus $g$ (cf. SHMI). Theorem 10 also proves that the “Jacobi schemes of
level $n$” in the theory of moduli are nonsingular, and indeed smooth over
$\mathbf Z$.

\src{208}
Return to the situation and notation of Theorem 9, now supposing that $A$ is a
finite-dimensional local algebra over $k$, which for simplicity is assumed
algebraically closed. Then $F(A)$ may be regarded as a covariant functor of
$A$ with values in the category of sets: a homomorphism of $k$-algebras
$A\longrightarrow B$ defines a map $F(A)\longrightarrow F(B)$, since every
flat $A$-scheme $X$ such that $X\otimes_A k=X_0$ gives rise to a $B$-scheme
$X\otimes_A B$ with the same properties. Suppose that one can find a complete
Noetherian local $k$-algebra $\mathcal O$ and a functorial isomorphism
\begin{equation}
  \tag{*}
  \operatorname{Hom}(\mathcal O,A)\xrightarrow{\sim}F(A)
\end{equation}
(where the first member denotes homomorphisms of $k$-algebras). It is easy to
see that such an $\mathcal O$ is determined up to canonical isomorphism. The
formal spectrum $\mathfrak Y$ of $\mathcal O$ (i.e. the one-point topological
space equipped with the sheaf of topological rings reduced to $\mathcal O$)
is then called the \emph{formal moduli scheme} for $X_0$. (It need not exist.)
Let $\mathfrak r$ be the maximal ideal of $\mathcal O$, and for every $n$ put
$\mathcal O_n=\mathcal O/\mathfrak r^{n+1}$ (so $\mathcal O_0=k$). The
canonical homomorphism $\mathcal O\longrightarrow\mathcal O_n$ is an element
of $\operatorname{Hom}(\mathcal O,\mathcal O_n)$, and hence defines an element
of $F(\mathcal O_n)$, i.e. a flat $\mathcal O_n$-scheme $X_n$ whose reduction
modulo $\mathfrak r$ is $X_0$. These $X_n$ are obtained from one another by
extension of scalars (here, by reduction), and consequently arise from a
well-determined formal scheme $\mathfrak X$ over the formal moduli scheme
$\mathfrak Y$; $\mathfrak X$ is flat over $\mathfrak Y$ and
$\mathfrak X_0=X_0$. The isomorphism (*) is then given, as is immediately
seen, by associating to every homomorphism of $k$-algebras
$\mathcal O\longrightarrow A$ the class of the $A$-scheme
$\mathfrak X\otimes_{\mathcal O}A$ (i.e. to every morphism of $k$-schemes
$Y'=\operatorname{Spec}(A)\longrightarrow\mathfrak Y$, one associates the
$Y'$-scheme $\mathfrak X\otimes_{\mathfrak Y}Y'$ obtained by base change).
Moreover, the isomorphism (*) and its preceding description remain valid if
one assumes only that $A$ is a complete Noetherian local $k$-algebra (not
necessarily Artinian). Of course, as usual, $\mathcal O$ may very well have
nilpotent elements \textit{a priori}, and it seems probable that cases should
exist in which $\mathcal O$ is itself Artinian without being equal to $k$.
This shows how inadequate the Kodaira--Spencer point of view (which restricts
to $A$ that are regular rings) is \textit{a priori} in the general case.

\src{209}
It remains to give sufficient conditions for a formal moduli scheme to exist
for $X_0$, assumed proper over $k$. In general, it is easy to give simple
necessary and sufficient conditions on a functor $A\longmapsto F(A)$ (from
finite-dimensional local $k$-algebras to sets) for it to be expressible in the
form $\operatorname{Hom}(\mathcal O,A)$ for a suitable $\mathcal O$. We shall
not detail them here. Let us merely point out that, in the case at hand, these
conditions impose nontrivial cohomological conditions on $X_0$, and it seems
unlikely that they always hold, although the speaker has not constructed a
counterexample. On the other hand, it seems plausible that
$H^0(X_0,\Theta_{X_0/k})=0$ is a sufficient condition (though by no means a
necessary one) for the existence of a formal moduli scheme. We confine
ourselves here to stating a theorem in a particularly simple case (whose
analogue in the theory of analytic spaces is well known; cf.
Kodaira--Spencer), which follows without difficulty from the results of the
preceding section.

\result{THEOREM}{10} Let $X_0$ be a scheme proper and nonsingular over the
field $k$, such that
\[
  H^0(X_0,\Theta_{X_0/k})=H^2(X_0,\Theta_{X_0/k})=0.
\]
Then a formal moduli scheme exists for $X_0$, corresponding to a regular local
ring $\mathcal O$ (i.e. a formal power-series algebra over $k$).

As already noted, it is not true in general that the formal scheme
$\mathfrak X$ over $\mathcal O$ is algebraizable; it is nevertheless known to
be so when $X_0$ is projective and $H^2(X_0,\Osh_{X_0})=0$ (Theorem 4), for
example when $X_0$ has dimension 1. This lends some support to the hope of
constructing a moduli scheme over the integers for curves of a given
genus\ldots{}

Methods such as those presented in this section can also be applied to the
construction and study of Picard varieties, as well as to many other
constructions. We shall return to this soon.

\fsection{8}{Application to the fundamental group}

The techniques presented above make it possible to undertake a systematic
study of the fundamental group, modeled on topological theory. The first two
theorems stated in this section generalize results of a recent paper by
Lang--Serre.

Let $X$ be a scheme. An $X$-scheme $X'$ is called an unramified covering
of $X$ if
\begin{enumerate}
  \item[\textit{i.}] $X'$ is finite over $X$, i.e. it is defined by a coherent
    sheaf of algebras $\mathcal A=\mathcal A(X')$ on $X$;
  \item[\textit{ii.}] $\mathcal A$ is a locally free sheaf on $X$;
  \item[\textit{iii.}] for every $x\in X$, the quotient
    \[
      \mathcal A_x/\mathfrak m_x\mathcal A_x
      =\mathcal A_x\otimes_{\Osh_x}\kappa(x)
    \]
    is a separable algebra over $\kappa(x)$.
\end{enumerate}

\src{210}
This notion of unramified covering (due to Serre and the speaker) has all the
elementary properties one can reasonably expect, which we shall not list. Let
us merely say that it gives rise to a Galois theory modeled on classical
Galois theory (and containing it, with proofs that are rather simpler than the
usual proofs of the latter) and on the Galois theory of topological coverings.
More precisely, call a \emph{geometric point} of a scheme $X$ a morphism $a$
from the spectrum $\xi$ of an algebraically closed field $\Omega$ into $X$,
i.e. the datum of an algebraically closed extension of the residue field
$\kappa(x)$ of a point $x=|a|$ of $X$ (called the \emph{locality} of the
geometric point $a$). If $X'$ is an unramified covering of $X$, associate to it
the set $E_a(X')$ of “geometric points of $X'$ above $a$,” i.e. the set of
pairs consisting of a point $x'\in X'$ above $x$ and a
$\kappa(x)$-homomorphism $\kappa(x')\longrightarrow\Omega$. Thus, for fixed
$(X,a)$, one obtains a functor $F(X,a)$ from the category $R(X)$ of unramified
coverings $X'$ of $X$ to the category of finite sets. If $X$ is connected, the
pair formed by $R(X)$ and $F(X,a)$ has the formal properties needed to be
isomorphic to the analogous pair defined by a suitable compact totally
disconnected topological group $\pi$ (i.e. an inverse limit of finite groups):
one takes the category $C(\pi)$ of finite sets $E$ on which $\pi$ acts
continuously, together with the forgetful functor $F(\pi)(E)=E$ from this
category to the category of finite sets. Moreover, $\pi$ is determined up to
canonical isomorphism by the condition that $(C(\pi),F(\pi))$ be isomorphic to
a given pair. Here $\pi$ is called the \emph{fundamental group} of the
connected scheme $X$ at the geometric point $a$, and is denoted $\pi_1(X,a)$.
If $X$ is not connected, it is replaced by the connected component of
$x=|a|$. If $X$ is connected, however, then the groups $\pi_1(X,a)$ for two
geometric points $a',a''$ of $X$ are isomorphic (the isomorphism being
determined only up to inner automorphisms), so that, as usual, one may choose
$a$ to suit the problem, for example at the generic point of
$X$, assumed irreducible. Of course, $\pi_1(X,a)$ is a covariant functor of the
pointed scheme $(X,a)$. Every statement concerning the classification of
unramified coverings\footnote{The printed text says “inseparable coverings.”
The paragraph concerns the category $R(X)$ of unramified coverings and its
equivalence with finite sets carrying a continuous $\pi_1(X,a)$-action.} can
then be translated into group-theoretic language according to the familiar
dictionary (except that here one must take account of the fact that the groups
are topological).

\src{211}
Our aim is to develop the analogue of the homotopy exact sequence for fiber
spaces, relative to a proper morphism $f:X\longrightarrow Y$. Since we do not
know what the higher homotopy groups should be, the results will necessarily
be incomplete. To apply the fundamental theorems of Section 3, we must first
make explicit a few elementary lemmas concerning schemes over fields or
Artinian rings (in accordance with the general procedure!).

\result{LEMMA}{1} Let $(X',a')$ be a pointed unramified covering associated to
a pointed representation of $\pi_1(X,a)$ on a finite set $E$ (with marked
point $e$). Then the canonical morphism
$\pi_1(X',a')\longrightarrow\pi_1(X,a)$ identifies the former group with the
stabilizer of $e$ in $\pi_1(X,a)$ (and is therefore injective).

\result{LEMMA}{2} Let $X$ be an algebraic scheme over a field $k$, and let
$k'$ be a radicial extension of $k$. Then every unramified covering of
$X\otimes_k k'$ is obtained by inverse image (i.e. extension of scalars) from
an unramified covering of $X$, determined up to canonical isomorphism.

It follows in particular from these two lemmas that for every algebraic
extension $K$ of $k$, and every geometric point $a'$ of
$X'=X\otimes_k K$ mapping to the geometric point $a$ of $X$, the functorial
homomorphism $\pi_1(X',a')\longrightarrow\pi_1(X,a)$ is injective.

\result{LEMMA}{3} Let $X$ be a complete scheme over a local Artinian ring $A$
such that $H^0(X,\Osh_X)=A$. Let $X'$ be an unramified covering of $X$, and
put $A'=H^0(X',\Osh_{X'})$, which is therefore a finite ring over $A$ (and may
\textit{a priori} be ramified over $A$). Let $X_0,X'_0$ be the associated
reduced subschemes of $X,X'$ (obtained by quotienting by the sheaves of
nilpotent elements in $\Osh_X,\Osh_{X'}$), and let $k$ be a subfield of
$A/\mathfrak r(A)$ over which $A/\mathfrak r(A)$ is finite (so $X_0$ is a
complete algebraic scheme over $k$, and $X'_0$ is an unramified covering of
it). Finally, let $\Omega$ be an algebraically closed extension of $k$, and
consider the unramified covering $X'_0\otimes_k\Omega$ of
$X_0\otimes_k\Omega$. Then the following two conditions are equivalent:
\begin{enumerate}
  \item[\textit{i.}] $X'_0\otimes_k\Omega$ is completely split over
    $X_0\otimes_k\Omega$.
  \item[\textit{ii.}] The natural morphism
    $X'\longrightarrow X\otimes_A A'$ is an isomorphism.
\end{enumerate}
Under these conditions, $A'$ is an unramified extension of $A$. Finally, if
$X'$ is connected, condition \textit{(i)} is equivalent to the following
apparently weaker condition:
\begin{enumerate}
  \item[\textit{i bis.}] $X'_0\otimes_k\Omega$ admits a regular section over
    $X_0\otimes_k\Omega$.
\end{enumerate}
When condition \textit{(ii)} holds, we shall say that the unramified covering
$X'$ of $X$ is \emph{geometrically trivial}.

\src{212}
\result{LEMMA}{4} Let $f:X\longrightarrow Y$ be a proper morphism such that
$f_*(\Osh_X)=\Osh_Y$. Let $a$ be a geometric point of $X$, and let $b$ be its
image in $Y$. Then $\pi_1(X,a)\longrightarrow\pi_1(Y,b)$ is surjective.

Indeed, it is enough to prove the following: if an unramified covering $Y'$ of
$Y$ (corresponding to a locally free sheaf of algebras $\mathcal A$) is such
that $X\otimes_Y Y'$ is disconnected, then $Y'$ is disconnected as well.
Indeed, $f^*(\mathcal A)$ is then a direct sum of two nonzero sheaves of rings,
and hence so is its direct image; but the latter is precisely
$\mathcal A\otimes f_*(\Osh_X)=\mathcal A$.

\result{LEMMA}{5} Let $X$ be a complete scheme over a field $k$. Suppose that
$H^0(X,\Osh_X)$ is a local ring $A$ and that $A/\mathfrak r(A)$ is radicial
over $k$. Let $\Omega$ be an algebraic closure of $k$, and put
$\overline X=X\otimes_k\Omega$ (which is connected). Choose a geometric point
$\overline a$ of $\overline X$ mapping to the geometric point $a$ of $X$.
Then there is an exact sequence
\[
  e\longrightarrow\pi_1(\overline X,\overline a)
  \longrightarrow\pi_1(X,a)
  \longrightarrow\pi_1(k,b)\longrightarrow e
\]
(where $\pi_1(k,b)$ is the Galois group of $\Omega$ over $k$).

The injectivity of the first homomorphism was seen in Lemmas 1 and 2;
exactness in the middle follows from Lemma 3; and the surjectivity of the last
homomorphism (the only part that uses the fact that $A/\mathfrak r(A)$ is
radicial) follows from Lemma 4.

\result{PROPOSITION}{4} Let $f:X\longrightarrow Y$ be a proper flat morphism
such that for every $y\in Y$ the algebra
$H^0(f^{-1}(y),\Osh_{f^{-1}(y)})$ is separable over the residue field
$\kappa(y)$ (as is the case, for example, if $f^{-1}(y)$ is a separable scheme
over $\kappa(y)$, i.e. reduced and such that the fields corresponding to its
irreducible components are separable extensions of $\kappa(y)$). Then the
covering
$Y'$ of $Y$ associated with $f_*(\Osh_X)$ is unramified.

The proof is easy from Theorem 2.

\src{213}
Together with Lemma 1, this proposition reduces the homotopy-theoretic study
of proper flat morphisms (with separable fibers), for all practical purposes,
to the case $f_*(\Osh_X)=\Osh_Y$ (for, using Stein factorization, one replaces
$Y$ by $Y'$).

\result{REMARK}{} A flat morphism of finite type whose fibers are separable
(respectively nonsingular) schemes is called \emph{separable} (respectively
\emph{smooth}). One proves that if $f$ is flat and $f^{-1}(y)$ is separable
(respectively nonsingular), then there is a neighborhood of $f^{-1}(y)$ on
which $f$ is separable (respectively smooth). The same result holds for
“absolutely normal” (“Bertini's theorem”).

Let $f:X\longrightarrow Y$ be a proper morphism such that
\begin{equation}
  \tag{i}
  f_*(\Osh_X)=\Osh_Y,
\end{equation}
and let $X'$ be a scheme finite over $X$. Let $Y'$ be the covering of $Y$
corresponding to the Stein factorization of $X'\longrightarrow Y$ (cf.
Theorem 5). Let $y\in Y$. The set of connected components of the fiber $F'$ of
$X'$ over $y$ therefore identifies with the set of points $y'$ of $Y'$ above
$y$ (Theorem 5). Consider the evident morphism
\begin{equation}
  \tag{*}
  X'\longrightarrow X\times_Y Y'
\end{equation}
induced by the natural morphisms $X'\longrightarrow X$ and
$X'\longrightarrow Y'$.\footnote{The printed text omits the prime on the
second $X'$; the morphism to $Y'$ is the one in the Stein factorization of
$X'\to Y$.} This is an isomorphism whenever $X'$ has the form
$X\times_Y Y''$, where $Y''$ is an unramified covering of $Y$; then $Y'$ is
simply $Y''$ and (*) is the identity. We seek conditions under which $X'$ has
the form just described, i.e. under which $Y'$ is unramified and (*) is an
isomorphism. To this end, let $F$ be the fiber of $X$ over $y$; it is a proper
scheme over $\kappa(y)$, and $F'$ is a covering of it (unramified if $X'$ is).
Let $F'_1$ be a connected component of $F'$ corresponding to a point $y'_1$
of $Y'$ above $y$. Suppose in addition that
\begin{enumerate}
  \item[\textit{(ii)}] $X'$ is unramified over $X$ at the points of $F'_1$
    (hence $F'_1$ is an unramified covering of $F$), and
  \item[\textit{(iii)}] $F'_1$ is a geometrically trivial covering of $F$
    (cf. Lemma 3).
\end{enumerate}

\result{THEOREM}{11} Under these conditions, there exists an open neighborhood
$U'$ of $y'_1$ in $Y'$ such that (*) is an isomorphism above $U'$. Moreover,
$Y'$ is unramified at $y'_1$ over $Y$ (although it may be ramified at the other
points $y'$ of $Y'$ above $y$).

The conditions \textit{(ii)} and \textit{(iii)} are of course also necessary
for the conclusion of the theorem to hold. The proof is easy, using Lemma 3
and Theorem 2.

\src{214}
\result{COROLLARY}{1} Continue to assume \textit{(i)}. An unramified covering
$X'$ of $X$ is isomorphic to the inverse image of an unramified covering $Y'$
of $Y$ if and only if $X'$ induces a geometrically trivial covering on every
fiber $f^{-1}(y)$.

By Theorem 11, the set of points of $Y$ for which this condition holds is open,
so it is enough to verify it at the closed points $y$\ldots{} The following
statement is equivalent to Corollary 1: the kernel of the homomorphism
$\pi_1(X)\longrightarrow\pi_1(Y)$ (surjective by Lemma 4) is the closed normal
subgroup generated by the images in $\pi_1(X)$ of the groups
$\pi_1(\overline{f^{-1}(y)})$, where
\[
  \overline{f^{-1}(y)}
  =f^{-1}(y)\otimes_{\kappa(y)}\overline{\kappa(y)},
\]
and $\overline{\kappa(y)}$ denotes an algebraic closure of $\kappa(y)$. Since
one cannot choose the same base point for all the fibers, the homomorphisms
$\pi_1(\overline{f^{-1}(y)})\longrightarrow\pi_1(X)$ are in any event
determined (once a base point for $X$, and hence for $Y$, has been chosen) only
up to composition with an inner automorphism of $\pi_1(X)$.

\result{COROLLARY}{2} Under the general hypotheses of Theorem 11, suppose in
addition that $Y,X,X'$ are integral, and let $K,L,L'$ be their function fields.
Then there exists a separable intermediate extension $K'$ of $K$ in $L'$,
linearly disjoint from $L$, such that $L'=LK'$ (whence
$L'=L\otimes_K K'$).

(Apply the last part of Lemma 3 to the geometric fiber of $X$.) The most
interesting application of Theorem 11 occurs when $f$ is separable. Then $X'$
is also separable over $Y$, so by Proposition 4, $Y'$ is unramified over $Y$;
hence the second member $X\times_Y Y'$ in (*) is unramified over $X$. It
follows easily that:

\result{COROLLARY}{3} Suppose, in addition to \textit{(i)}, that $f$ is
separable. Let $X'$ be a connected unramified covering of $X$. Then $X'$\footnote{The
printed text has “$X$” here; the subject of the criterion is the covering $X'$
introduced in the preceding sentence.} is
the inverse image of an unramified covering $Y'$ of $Y$ if and only if the
covering $\overline F'$ induced on one geometric fiber
$\overline F=f^{-1}(y)$ admits a regular section.

Note that it was not necessary to suppose that $\overline F'$ was
geometrically trivial over $\overline F$ (this follows \textit{a posteriori},
although \textit{a priori} it is a much stronger condition). Corollary 3 is in
fact equivalent to the following statement.

\src{215}
\result{COROLLARY}{4} Let $f:X\longrightarrow Y$ be a proper separable
morphism such that $f_*(\Osh_X)=\Osh_Y$. Let $\overline F$ be the geometric
fiber of a point $y\in Y$, and choose a geometric point in $\overline F$; the
morphisms $\overline F\longrightarrow X\longrightarrow Y$ then give geometric
points in $X,Y$, which are taken as base points for the fundamental groups of
$\overline F,X,Y$. Under these conditions there is an exact sequence
\[
  \boxed{
    \pi_1(\overline F)\longrightarrow\pi_1(X)
    \longrightarrow\pi_1(Y)\longrightarrow e
  }.
\]

The following two statements of Serre--Lang, now freed of any normality
hypothesis, follow easily.

\result{COROLLARY}{5} Let $X,Y$ be connected schemes over a field $k$. Suppose
that the reduced scheme $X_{\mathrm{red}}$ is separable over $k$ (as is
automatically the case if $k$ is perfect), and that $X$ or $Y$ is proper over
$k$.\footnote{Condition added by the 1962 erratum.} Choose a geometric
point $a$ (respectively $b$) in $X$ (respectively $Y$), hence a geometric point
$c=(a,b)$ in $X\times_kY$ and a natural morphism
\[
  \pi_1(X\times_kY,c)\longrightarrow
  \pi_1(X,a)\times\pi_1(Y,b)
\]
(induced by the functorial morphisms from $\pi_1(X\times Y,c)$ to
$\pi_1(X,a)$ and $\pi_1(Y,b)$). This morphism is injective, and is even
bijective if $k$ is algebraically closed.

(Surjectivity in the latter case is more or less trivial.) One deduces, with
Serre--Lang:

\result{COROLLARY}{6} Let $X$ be a connected proper algebraic scheme over an algebraically closed field $k$, and let $K$ be an algebraically closed
extension of $k$. Then the fundamental groups of $X$ and $X\otimes_kK$ are
the same; i.e. every unramified covering of the latter scheme is obtained by
extension of scalars from an unramified covering of $X$, unique up to
isomorphism.

\result{REMARKS}{}

1) Using Proposition 4, one sees that the hypothesis
$f_*(\Osh_X)=\Osh_Y$ in Corollary 4 is not essential. In the general case,
instead of placing the unit group $e$ after $\pi_1(Y)$, one should continue
with
\[
  \pi_0(\overline F)\longrightarrow\pi_0(X)
  \longrightarrow\pi_0(Y)\longrightarrow e.
\]
as in algebraic topology.

\src{216}
2) In general, nothing can yet be said about the kernel of
$\pi_1(\overline F)\longrightarrow\pi_1(X)$, which ought to involve a
$\pi_2(Y)$. It nevertheless seems that one should be able to prove that
$\pi_1(\overline F)\longrightarrow\pi_1(X)$ is injective when $Y$ is the
spectrum of a local ring $A$, using Theorem 12 below (which asserts this when
$A$ is complete).

Theorem 11 used only Theorems 1 and 2. We now apply Theorem 3, using the
following elementary lemma.

\result{LEMMA}{6} Let $X$ be a scheme and let $X_0$ be the corresponding
reduced scheme (i.e. with the nilpotent elements killed). Then every
unramified covering $X'_0$ of $X_0$ is induced by an unramified covering $X'$
of $X$, determined up to canonical isomorphism.

This purely local lemma plays a role here analogous to that of Theorem 8 in
the theory of moduli. Combining it with the existence theorem (Theorem 3)
gives:

\result{THEOREM}{12} Let $A$ be a complete Noetherian local ring with residue
field $k$, and let $X$ be a scheme proper over $A$. Then every unramified
covering $X'_0$ of $X_0=X\otimes_A k$ is induced by an unramified covering
$X'$ of $X$, unique up to isomorphism.

In other words:

\result{COROLLARY}{1} Choose a geometric point in $X_0$ as base point for the
fundamental groups of $X_0$ and $X$. Then the canonical homomorphism
$\pi_1(X_0)\longrightarrow\pi_1(X)$ is an isomorphism.

Now apply Lemma 5 to $X_0$ (assuming for simplicity that
$H^0(X_0,\Osh_{X_0})=k$), and observe that since $A$ is complete, its
unramified extensions correspond to the unramified extensions of its residue
field, i.e. $\pi_1(Y)=\pi_1(k)$ (where $Y=\operatorname{Spec}(A)$). This gives
the exact sequence
\[
  e\longrightarrow\pi_1(\overline X_0)
  \longrightarrow\pi_1(X)\longrightarrow\pi_1(Y)\longrightarrow e.
\]

\result{COROLLARY}{2} Let $f:X\longrightarrow Y$ be a proper flat morphism,
let $y_1$ be a point of $Y$, and let $y_0$ be a specialization of $y_1$.
Consider the corresponding “geometric” fibers $\overline X_1$ and
$\overline X_0$, and suppose that $\overline X_0$ is separable and connected
(which implies that $\overline X_1$ satisfies the same conditions). Then one
can find a group homomorphism
$\pi_1(\overline X_1)\longrightarrow\pi_1(\overline X_0)$, defined up to an
inner automorphism, and this homomorphism is surjective.

\src{217}
One might hope that this homomorphism is always bijective. Unfortunately this
is not so in general when $\kappa(y_0)$ has characteristic $>0$. Below we shall
nevertheless obtain a bound on the kernel of this homomorphism (at least when
$\overline X_0$ is nonsingular), implying that if $\kappa(y_0)$ has
characteristic $0$, then the homomorphism above is bijective (a result that
could also be proved by transcendental methods). In any case, we already
obtain a bound on the $\pi_1$ of a special fiber in terms of that of a generic
fiber. For example, using the fact that an algebraic curve in characteristic
$p$ lifts to a curve in characteristic $0$ (Theorem 9, Corollary 4), one finds
by transcendental methods:

\result{COROLLARY}{3} Let $X_0$ be the scheme of a complete nonsingular curve
over an algebraically closed field of arbitrary characteristic, and let $g$ be
the genus of $X_0$. Then $\pi_1(X_0)$ admits $2g$ topological generators,
subject to the familiar relation.

By a familiar hyperplane-section technique one deduces:

\result{COROLLARY}{4} Let $X$ be a nonsingular projective scheme over an
algebraically closed field of arbitrary characteristic. Then $\pi_1(X)$
admits a finite number of topological generators.

Let us seek the kernel of the homomorphism
$\pi_1(\overline X_1)\longrightarrow\pi_1(\overline X_0)$. For this purpose we
may suppose that $Y$ is the spectrum of a complete discrete valuation ring
$V=\Osh_y$ ($y=y_0$). The question is equivalent to the following: given an
unramified covering $\overline X'_1$ of $\overline X_1$ (which may, if desired,
be assumed Galois), determine under what conditions it comes from an
unramified covering of $X_0$. \textit{A priori}, the given covering is obtained
by extension of scalars from an unramified covering $X'_1$ of
$X_1\otimes_K K'$, where $K'$ is a finite subextension of the algebraic closure
$\overline K$ of the fraction field $K$ of $V$; if $\overline X'_1$ is Galois
with group $G$, one may choose $X'_1$ Galois with group $G$. One then has: for
$X'_1\otimes_{K'}\overline K=\overline X'_1$ to come from an unramified
covering of $X_0$, it is necessary and sufficient that there exist a finite
extension $K''$ of $K'$ in $\overline K$ such that
$X''_1=X'_1\otimes_{K'}K''$ has the form $X''\otimes_{V''}K''$, where $V''$ is
the integral closure of $V$ in $K''$ and $X''$ is an unramified covering of
$X\otimes_VV''$. Suppose, for example, that $X_0$ is absolutely normal. It
follows that $X\otimes_VV''$ is normal (being flat over $V''$ with normal
special fiber), and its function field is the same field $K''(X_1)$ as that of
\[
  X_1\otimes_KK''
  =(X\otimes_VK)\otimes_KK''
  =X\otimes_VK''
  =(X\otimes_VV'')\otimes_{V''}K''.
\]
Let $L''=K''(X'_1)$ be the function field of
$X'_1\otimes_{K'}K''$. It is a finite separable extension of $K''(X_1)$, and
the preceding condition also means that $L''$ is an unramified extension of
the function field of $X\otimes_VV''$ (i.e. the normalization of
$X\otimes_VV''$ in $L''$ is unramified over $X\otimes_VV''$). It is enough to
check that $L''$ is unramified at the points of the special fiber of
$X\otimes_VV''$ (since it is unramified over the generic fiber
$X_1\otimes_KK''$). If $X_0$ is nonsingular, the “purity theorem” of
Nagata--Zariski shows that it is enough to check that $L''$ is unramified over
the local ring $O''$ at the generic point of the special fiber of
$X\otimes_VV''$. This is a discrete valuation ring, the normalization in
$K''(X_1)$ of the local ring $O\subset K(X_1)$ at the generic point of the
special fiber of $X$. Thus we have obtained:

\src{218}
\result{COROLLARY}{5} Under the preceding hypotheses and notation, the
unramified covering $X'_1\otimes_{K'}\overline K$ of
$\overline X_1=X_1\otimes_K\overline K$ comes from an unramified covering of
$\overline X_0$ if and only if there exists a finite subextension $K''$ of
$\overline K/K'$ such that $K''(X'_1)$ is unramified over the discrete
valuation ring $O''\subset K''(X_1)$.

Now observe that $O''$ is the normalization in $K''(X_1)$ of the discrete
valuation ring $O'\subset K'(X_1)$ (the normalization of $O$ in $K'(X_1)$),
and that $O'$ contains the normalization $V'$ of $V$ in $K'$, a uniformizer
$u$ of $V'$ being also a uniformizer of $O'$. Suppose now that $X'_1$ is
Galois with Galois group $G$ of order $n$ prime to the characteristic $p$ of
$\kappa(y_0)$ (which is also the characteristic of the residue field of
$O'$). Thus $K'(X'_1)$ is tamely ramified over $O'$, and it follows easily
(“Abhyankar's lemma”) that after adjoining an $n$th root $v$ of a uniformizer
of $O'$, it becomes unramified over the normalization of $O'$ in
$K'(X_1)(v)$. But one may take for $v$ an $n$th root of a uniformizer of $V'$,
which proves that the condition of Corollary 5 is satisfied. (The possibility
of using Abhyankar's lemma and the purity theorem was sold to me by Serre.) To
express the result obtained, for every compact totally disconnected group
$\pi$ let $\widetilde\pi$ denote the quotient of $\pi$ by the closed subgroup
generated by its Sylow $p$-subgroups, i.e. the inverse limit of the discrete
quotient groups of $\pi$ whose orders are prime to $p$. With this notation one
obtains:

\result{THEOREM}{13} Let $f:X\longrightarrow Y$ be a proper flat morphism,
let $y_1$ be a point of $Y$, and let $y_0$ be a specialization of $y_1$.
Suppose that $\overline X_0$ is connected and nonsingular. Then the
homomorphism
$\widetilde\pi_1(\overline X_1)\longrightarrow
\widetilde\pi_1(\overline X_0)$ induced by the homomorphism that is surjective
by Theorem 12, Corollary 2, is an isomorphism.

In other words:

\src{219}
\result{COROLLARY}{1} The classification of unramified Galois coverings whose
Galois group has order prime to the characteristic $p$ of $\kappa(y_0)$ is the
same for $\overline X_0$ and $\overline X_1$.

In particular, if $\kappa(y_0)$ has characteristic zero, one obtains
algebraically that
$\pi_1(\overline X_1)\longrightarrow\pi_1(\overline X_0)$ is bijective.

Finally, the techniques used also give the following result, which is more
general than Theorem 13.

\result{THEOREM}{14} Let $f:X\longrightarrow Y$ be a proper smooth morphism,
and let $D$ be a closed subscheme of $X$, smooth over $Y$ and of codimension 1
at all its points. For a fiber $Z=f^{-1}(z)$ of $f$, put $Z'=Z-Z\cap D$ and
let $\pi_1^t(Z')$ be the quotient of the fundamental group $\pi_1(Z')$ that
classifies the unramified coverings of $Z'$ which are tamely ramified above
$\overline{Z\cap D}$. Let $y_0,y_1$ be as in Theorem 13. Then there is a
surjective homomorphism (defined up to an inner automorphism)
\[
  \pi_1^t(\overline X'_1)\longrightarrow\pi_1^t(\overline X'_0),
\]
and the corresponding homomorphism
\[
  \widetilde\pi_1^t(\overline X'_1)
  \longrightarrow\widetilde\pi_1^t(\overline X'_0)
\]
is an isomorphism.

Corresponding variants of the corollaries to Theorem 13 and of Corollary 4 to
Theorem 12 follow. Likewise, using Theorem 9, Corollary 3, one finds by
transcendental methods:

\result{COROLLARY}{} Let $X_0$ be the scheme of a complete nonsingular curve
over an algebraically closed field of arbitrary characteristic, and let
$S=(s_i)_{1\le i\le n}$ be a finite subset of $X_0$ with $n$ elements. Then
$\pi_1^t(X_0-S)$ admits $2g+n$ topological generators
$x_i,y_i$ ($1\le i\le g$) and $\sigma_j$ ($1\le j\le n$), subject to the
relation
\[
  \left(\prod_i x_i y_i x_i^{-1}y_i^{-1}\right)
  \sigma_1\cdots\sigma_n=1,
\]
where the $\sigma_j$ are generators of the inertia groups corresponding to the
$s_j$. For every finite group $G$ of order prime to the characteristic,
generated by elements $x_i,y_i,\sigma_j$ satisfying the preceding relation,
there exists an unramified Galois covering of $X_0-S$ with group $G$, whose
inertia groups at the points $s_j$ are generated by the $\sigma_j$.

When $X_0$ has genus 0 and $n=3$, this gives a solution of the “Problem
of the three points,” at least for Galois coverings of order prime to the
characteristic. (Here Theorem 9 is unnecessary, and it seems moreover that the
preceding corollary can be deduced from the special case considered in the
problem of the three points.)

\src{220}
\result{REMARKS}{}

1) A more complete study, no doubt involving generalized Galois coverings of
$X,X_0,X_1$ (with possibly infinitesimal Galois group), should make it possible
to recover the kernel in Theorem 12, Corollary 2. By contrast, a study of
coverings admitting ramification that is not tame seems likely to be much more
difficult.

2) Lemma 6, together with a result of Grauert concerning the formal completion
of a nonsingular projective scheme along a hyperplane section (or with the
theorem, not yet proved, mentioned in Remark 2 after Theorem 11), also makes it
possible to prove in “abstract” algebraic geometry the classical Lefschetz
theorem on the fundamental group.

\begin{center}
  \textbf{BIBLIOGRAPHY}
\end{center}

\begin{description}
  \item[{[1]}] Dieudonné (J.) and Grothendieck (A.). --- \emph{Éléments de
    géométrie algébrique}, Publications mathématiques de l'Institut des Hautes
    Études Scientifiques (forthcoming).
  \item[{[2]}] Grothendieck (Alexander). --- \emph{The cohomology theory of
    abstract algebraic varieties}. International Congress of Mathematicians
    [1958. Edinburgh] (forthcoming).
  \item[{[3]}] Serre (Jean-Pierre). --- \emph{Faisceaux algébriques
    cohérents}, Annals of Math., Series 2, vol. 61, 1955, pp. 197--278.
  \item[{[4]}] Serre (Jean-Pierre). --- \emph{Géométrie algébrique et
    géométrie analytique}, Ann. Institut Fourier Grenoble, vol. 6, 1955--1956,
    pp. 1--42.
\end{description}

\par\smallskip
\begingroup\footnotesize
\endgroup
